\chapter{Why dense attention is unavoidably $\bigO(n^2)$}
\label{ch:why-quadratic}

\newthought{Self-attention} from Chapter~\ref{ch:self-attention} forms an
\(n \times n\) matrix per head per layer. This chapter is the careful
accounting of what that costs in compute, in memory, and in dollars at the
context lengths people actually want.

\section{The compute count}

For one attention head with \(n\) tokens and head dimension \(\dhead\):

\begin{itemize}
\item Forming \(Q, K, V\) costs \(3 \cdot n \cdot d \cdot \dhead\) FLOPs
(three projection matmuls). Linear in \(n\).
\item Forming \(QK\T\) costs \(n^2 \cdot \dhead\) FLOPs. \emph{Quadratic in \(n\).}
\item The softmax costs \(\bigO(n^2)\) operations.
\item Forming \(AV\) costs \(n^2 \cdot \dhead\) FLOPs. Quadratic again.
\item The output projection costs \(n \cdot d^2\). Linear.
\end{itemize}
Summing across \(\heads\) heads and \(L\) layers, the dominant term is
\(\bigO(L \cdot \heads \cdot n^2 \cdot \dhead) = \bigO(L \cdot n^2 \cdot d)\).

\keyidea{Attention's compute scales as \(\bigO(n^2 d)\) per layer. Doubling
\(n\) does not double the work — it \emph{quadruples} it. Going from 32K to
1M tokens (a factor of \(\sim 32\)) makes attention \(\sim 1024\)$\times$
more expensive per layer.}

\section{The memory count}

Two kinds of memory bills come due.

\paragraph{Activation memory during a forward pass.} The attention matrix
\(A\) has \(n^2\) entries per head per layer. For \(n = 10^6\), one head,
fp16: \(2 \cdot 10^{12}\) bytes \(= 2\) TB.\sidenote{Per layer. Per head.
Even with bf16 or fp8 we are still in the multi-gigabyte range per head per
layer at \(n = 10^6\). This is why FlashAttention, which never
\emph{materialises} \(A\) in main memory, was such a big deal —
Chapter~\ref{ch:flashattention}.}

\paragraph{KV cache memory at inference.} During autoregressive generation,
to answer ``what comes after token \(t\)'' the model needs the keys and
values for tokens \(1, \ldots, t\). It would be wasteful to recompute them at
every step, so they are cached. The cache size grows linearly in \(t\):
\begin{equation}
  \text{KV bytes per token} \;=\; 2 \cdot L \cdot \heads \cdot \dhead \cdot s,
\end{equation}
where \(s\) is the bytes per scalar (2 for fp16). For a 7B-class model with
\(L=32\), \(\heads=32\), \(\dhead=128\), fp16: ${\sim}524{,}288$ bytes
\(\approx\) 0.5 MB \emph{per token}.\sidenote{The arithmetic: \(2 \cdot 32
\cdot 32 \cdot 128 \cdot 2 = 524{,}288\). One H100 has 80 GB of HBM. Divide:
${\sim}160{,}000$ tokens before the KV cache fills the entire card.
Frontier 70B-class models are roughly 2--4$\times$ worse per token.}

\paragraph{1M tokens hits the wall.} 0.5 MB/token \(\times\) \(10^6\) tokens =
524 GB of KV cache. The biggest single GPU available in 2026 is roughly
192 GB. A 1M-token transformer cannot keep its own working memory on a single
device, full stop. This is the engineering motivation for everything in
Part~IV (FlashAttention, MQA, GQA, paged caches) and the architectural
motivation for everything in Part~V (state-space models, sparse attention,
hybrids).

\section{Translated into dollars}

Run a \(70\)B model with dense attention at \(n = 10^6\) on a cluster of
H100s. Even with FlashAttention-3 reducing constants, an order-of-magnitude
estimate puts a single forward pass on the order of tens of GPU-seconds.
At cloud prices, that is dollars \emph{per inference}, before counting the
fact that the cache will not fit on one machine and so multi-GPU
communication adds another layer of slowdown.

This is what people mean when they say ``transformers don't scale to long
context''. The math works on paper. The hardware bill does not work on a
business model.

\keyidea{The engineering question for 2026 is not ``how do we make a model
with a 1M token context window'', it is ``how do we make a 1M token context
\emph{cheap enough to serve to actual users at scale}''.}

\section{Why this isn't the dimensions' fault}

Notice what we did \emph{not} count as quadratic: anything in \(d\). The
hidden dimension \(d\) shows up linearly. The vocabulary \(V\) shows up
linearly in the embedding/output layers. The number of layers \(L\) shows
up linearly. \emph{Only} sequence length \(n\) is quadratic.

That asymmetry is structural. Each token has to look at every other token
because attention is, by definition, an all-pairs operation. If you want to
remove the \(n^2\), you have to change \emph{which pairs} get scored. That
is the topic of Part~V.

\section{Bridge to scaling laws}

Before we discuss how to fix the \(n^2\), we have to address a different
question: even at \emph{short} context, how does the loss of a transformer
behave as we make it bigger? The 2020 paper that answered that question
\citep{kaplan2020scaling} is the reason any of this matters at all — it
turned LM training from research into engineering. Part~III is about that
paper and its 2022 correction.

\fillme{FILL-06-01}{Code listing: KV-cache size calculator}{%
  10-line Python function \texttt{kv\_cache\_bytes(n, L, H, d\_head, dtype\_bytes=2)}
  that returns the total KV-cache size in bytes for an MHA model. Then
  variants \texttt{kv\_cache\_bytes\_gqa(...)} and \texttt{kv\_cache\_bytes\_mqa(...)}.
  Print a 4-row table for the Llama-3 8B configuration at
  \(n \in \{8\text{K}, 128\text{K}, 1\text{M}, 10\text{M}\}\) with all three
  schemes. Caption: which row crosses 80~GB (one H100) first?}{%
  KV-cache equation in this chapter; Llama-3 config from FILL-04-02.}{%
  ${\sim}10$ lines of code + 4-row output table.}

\fillme{FILL-06-02}{Worked example: dollars per query at \(n = 128\text{k}\)}{%
  Order-of-magnitude estimate of the cost to do a single forward pass on a
  70B-parameter dense transformer at 128k tokens with FlashAttention-2 on
  H100s. Show: total FLOPs, per-GPU tokens/s ballpark, GPU-seconds, GPU-hours
  at on-demand cloud pricing (state assumption explicitly: \$3/GPU-hr).
  Repeat at \(n = 1\text{M}\) and note where the answer becomes
  uncomfortable. End with one sentence on why this is the entire reason
  Part~V exists.}{%
  Llama-3 70B paper for FLOP/token; H100 spec sheet (${\sim}2$ PFLOP/s
  bf16); current cloud pricing pages.}{%
  Half a page; one inline arithmetic chain; one short table.}
